\begin{enumerate}
\item Augmenter un prix de 20\,\% c'est multiplier ce prix par $\left(1 + \dfrac{20}{100}\right) = 1 + 0,20 = 1,2$.

Le nouveau prix est égale à : $400 \times 1,2 = 480$~(\euro).
\item Baisser un prix de 10\,\% c'est multiplier ce prix par $\left(1 - \dfrac{10}{100}\right) = 1 - 0,10 = 0,9$.

Le nouveau prix est égale à : $130 \times 0,9 = 117$~(\euro).
\item Le nouveau prix est $P \times 1,2 \times 1,2 = P \times 1,2^2$.
\item Ont voté pour A : $\dfrac 14 = \dfrac{1 \times 25}{4 \times 25} = \dfrac{25}{100}$.

Ont voté pour B : $\dfrac{20}{100}$ ;

Ont voté pour C : $\dfrac{1}{3} = \dfrac{1 \times 33,3}{3 \times 33,3} = \dfrac{33,3}{99,9} \approx \dfrac{33,3}{100}$ ;

Ont voté pour D à peu près $1 - \dfrac{25}{100} - \dfrac{20}{100} - \dfrac{33,3}{100} = 1 - \dfrac{78,3}{100} = \dfrac{21,7}{100}$.

Le plus faible pourcentage a été obtenu par le candidat B.
\item $A = \dfrac{2}{1 - \dfrac 23} = \dfrac{2}{\dfrac 13} = 2 \times \dfrac 31 = 2 \times 3 = 6$.
\item $A = \dfrac{1}{100} +  \dfrac{1}{\np{1000}} = 0,01 + 0,001 = 0,011$.
\item 75 minutes correspondent à 60 + 15 minutes soit 1 h et 15 minutes.

Or 15 (min) $= \dfrac{15}{60} = \dfrac{15 \times 1}{15 \times 4} = \dfrac 14 = 0,25$~(h).

Donc 75 minutes correspondent à 1,25~(h).
\item $10^{30}$ s'écrit 1 suivi de 30 zéros alors que  $10^{-30}$ s'écrit 0, suivis de 29 zéros et le chiffre 1 : ce nombre est négligeable par rapport au premier donc la somme est environ égale à  $10^{30}$.
\item Le coefficient directeur est négatif ce qui exclut $D_1$ et $D_2$ ;

L'ordonnée à l'origine (la valeur de $y$ quand $x = 0$) est positive ce qui exclut $D_4$. Il reste $D_3$.
\item $3x = 0$ si l'un des facteurs est nul : or $3 = 0$ ne donne pas de solution ; il reste donc $x = 0$.
\item $\dfrac{144}{x} = 9$ donne en multipliant par $x$ (supposé non nul sinon l'énoncé n'a pas de sens) $144 = 9x$ et en multipliant chaque membre par $\dfrac 19$ : $144 \times \dfrac 19 = x$, soit $x = \dfrac{144}{9}$.
\item La moyenne est égale à :

$\dfrac{1 \times 10 + 1 \times 13 + 1  \times 12 + 2 \times x}{1 + 1 + 1 + 2} = \dfrac{10 + 13 + 12 + 2x}{5} = \dfrac{35 + 2x}{5}$.

Si $\dfrac{35 + 2x}{5} = 15$ alors en multipliant par 5 : $35 + 2x = 75$, soit en ajoutant $- 35$,

 $2x = 40$ et enfin $x = 20$.
\end{enumerate}
